\section{Introduction}
\label{section:intro}

Converting standard five-line staff notation into numbered musical notation
(Jianpu) is a tedious yet essential task. In traditional Chinese music clubs and
ensembles, Jianpu is the standard for sight-reading and performance, and an
automated tool that directly outputs readable Jianpu would eliminate hours of
manual transcription, letting musicians focus on playing rather than on score
translation. Despite this practical demand, the problem is largely open: current
Optical Music Recognition (OMR) research overwhelmingly targets the digitisation
of \emph{Western} notation into intermediate formats such as MusicXML or MIDI.
Seminal pipelines~\citep{calvo2020understanding} and recent open-source
systems~\citep{oemer2023} rely on multi-stage semantic segmentation followed by
heuristic note-grouping rules; these stages accumulate cascading errors and are
not designed for cross-format (staff-to-Jianpu) translation.

We instead treat the task as a single image-to-sequence translation problem. A
Vision-Encoder-Decoder ingests a 2-D image of a printed monophonic score and
autoregressively decodes a 1-D sequence of custom semantic tokens, which a
post-processor maps to Jianpu. By avoiding intermediate representations and
hand-coded heuristics, the network must implicitly learn the music-theoretic
relationship between the 2-D spatial arrangement of symbols and the 1-D target.
Three challenges shape the design:
\begin{itemize}
    \item \textbf{Challenge 1: \textit{the 2D-to-1D semantic gap.}} Pitch is
    encoded by the \emph{absolute vertical position} of a notehead on the staff,
    while rhythm and ordering are largely horizontal. The model must reconcile
    both axes into one ordered token stream.
    \item \textbf{Challenge 2: \textit{translation without intermediate
    representations.}} Bypassing MusicXML/MIDI means the model learns alignment
    and notation conventions purely from visual supervision, with no explicit
    rule engine to fall back on.
    \item \textbf{Challenge 3: \textit{data scarcity.}} No large paired
    staff-to-Jianpu corpus exists, so training must rely on synthetic data while
    still generalising to real, photographed scores.
\end{itemize}

To make pitch and rhythm \emph{separately} measurable---and to expose where a
model fails---we decouple the target into four parallel streams
(\emph{type}, \emph{pitch}, \emph{rhythm}, \emph{attribute}) emitted in lock-step
by the generator, and give the decoder four corresponding output heads. This
design turns the architecture comparison into a controlled experiment: holding
the decoder, loss, and data fixed, we can attribute a pitch failure to the visual
representation rather than to the language model. Our main contributions are:

\begin{enumerate}
    \item \textbf{An end-to-end OMR$\rightarrow$Jianpu system.} A fully synthetic
    data pipeline (\texttt{music21}\,+\,\texttt{verovio}) that emits four decoupled
    label streams as ground truth, a multi-head Vision-Encoder-Decoder, a
    token-to-Jianpu post-processor, and a deployed interactive demo.
    \item \textbf{An architecture finding that reframes a ``ResNet cannot learn
    pitch'' result.} Under an autoregressive (AR) decoder, a pretrained ViT
    encoder solves the task (SER $0.0029$) while a from-scratch ResNet encoder
    structurally fails on pitch ($\approx0\%$). Attaching a BiLSTM\,+\,per-head CTC
    decoder to the \emph{same} ResNet features reaches a perfect SER of $0.0000$,
    showing the true bottleneck is the AR monotonic alignment over $1$-D column
    features---not the convolutional encoder.
    \item \textbf{A sample-efficiency study.} CRNN+CTC reaches usable transcription
    from $5{,}000$ samples whereas the ViT-AR pitch head does not begin to learn
    until $50{,}000$, contradicting the common expectation that ImageNet
    pretraining helps most in the low-data regime.
    \item \textbf{A sim-to-real study.} A photo-augmentation fine-tune more than
    halves SER on real phone photographs ($0.634\!\to\!0.258$), and a leave-one-out
    ablation isolates geometric (camera-angle) augmentation as the dominant factor.
\end{enumerate}

This is the Application track (NYCU 535354) final report; the code, synthetic
generator, trained checkpoints, and demo accompany the submission.
